\section{Introduction}
\label{sec:introduction}

\subsection{Cyclic proofs in proof assistants}

Cyclic proof systems offer a natural representation for reasoning about
inductive predicates: rather than packaging an explicit induction hypothesis
into every recursive case, the proof tree may contain \emph{back-edges} that
point to ancestor nodes, provided a global termination criterion---typically
the size-change principle of Lee, Jones, and Ben-Amram~\cite{lee2001size}---is
satisfied. Sprenger and Dam~\cite{sprenger2003structure} established the
equivalence between cyclic proofs and ordinary inductive proofs in the
\(\mu\)-calculus, formally justifying the soundness of cyclic-proof systems
for first-order predicates. Brotherston's PhD thesis~\cite{brotherston2006phd}
and the Cyclist prover~\cite{brotherston2012cyclist} demonstrated the
practical viability of cyclic proof search for sequent calculi with explicit
trace conditions.

Despite this theoretical foundation, cyclic proofs have not seen widespread
adoption in mainstream interactive proof assistants. The gap is partly
ergonomic: cyclic proofs are usually presented in a sequent-calculus surface
syntax with explicit bud-companion edges, which sits awkwardly inside the
tactic-mode workflows that dominate Lean, Coq, and Isabelle. It is also
partly methodological: a cyclic proof's soundness depends on a
non-local termination condition (size-change), and integrating this with the
proof assistant's kernel-level type-checking is a nontrivial engineering
question.

A separate strand of theoretical work has recently appeared in
Grotenhuis and Otten's \emph{Unravelling Abstract Cyclic Proofs into Proofs
by Induction}~\cite{grotenhuis2026unravelling}. They give an algorithm,
applicable to any sufficiently abstract cyclic system, that translates a
cyclic proof into an inductive proof by following the local annotation
data they call the \emph{Stack/Name/Reset} representation (§5) and the
\emph{relativised ancestor / inequality / hypothesis / cover} per-node data
(§6). Their Theorem~6.1 then justifies a direct structural-induction emission
based on this data. Crucially, their algorithm operates on abstract cyclic
representations, so the translation is independent of the specific
proof-system or term-language details.

\subsection{Contributions}

This thesis presents \emph{CyclicTactic}, a Lean~4 library that lets users
write cyclic proofs as ordinary Lean tactics and, by implementing the
Grotenhuis-Otten Theorem~6.1 unravelling, replaces the user's interactive
cyclic-style proof with a kernel-checked structural-induction theorem.

Concretely, the contributions are:

\begin{enumerate}
  \item \textbf{Side-channel tactic-mode integration.} Three primitive
    tactics---\code{cyclic <label>}, \code{cyc\_cases x with}, and
    \code{back <label> \{σ\}}---let the user write cyclic proofs inside
    Lean's standard tactic mode. Real Lean goals are visible at every
    cursor position; cyclic structure is recorded into a side-channel
    \code{ProofTree} value via source-position attribution
    (§\ref{sec:architecture}).

  \item \textbf{A paper-faithful implementation of Grotenhuis-Otten §5--§6.}
    We implement the Stack/Name/Reset annotation step of Definition~5.1, the
    Lemma~5.9 key-bounded fixpoint unfolding, the Theorem~5.2
    progressing-name verification, and the §6 per-node augmentation
    (\(\RelAnc\), \(\Ineq\), \(\Hyp\), \(\cov\), \(\sigma\)) needed to
    drive emission. To our knowledge this is the first implementation of
    the Grotenhuis-Otten algorithm in any proof assistant.

  \item \textbf{Theorem 6.1 emission with canonical-form replacement.}
    The §6-driven emitter (\code{Theorem6.Emit.translate}) produces a
    Lean tactic script. After the user's cyclic proof elaborates as a
    recursive \code{def} (giving real InfoView goals during writing), the
    environment is rolled back and the §6-emitted theorem is re-elaborated as
    the canonical user-facing declaration.

  \item \textbf{Mutual cyclic systems.} A \code{cyclic\_mutual} surface
    command, together with an event demultiplexer and a mutual emitter
    (\code{Theorem6.Emit.translateMutual}), handles cyclic proofs that span
    multiple mutually-recursive theorems with cross-companion back-edges.
    Cross-theorem buds emit as direct sibling-theorem calls, relying on
    Lean's mutual-recursion termination check for soundness.

  \item \textbf{Comparison with Cyclist on Brotherston's first-order
    benchmark.} On the cases that both systems handle, our emitter produces
    proofs that match the structure Cyclist would emit. On
    Brotherston's tests~10 and~11 (associativity and commutativity of
    \texttt{ADD}, where Cyclist fails because its first-order proof search
    cannot perform equational reasoning), we succeed by decoupling the
    cyclic structure (validated by the size-change check) from the per-arm
    reasoning (closed by Lean's tactic ecosystem).
\end{enumerate}

\subsection{Outline}

§\ref{sec:background} reviews cyclic proof theory, the size-change principle,
and the Grotenhuis-Otten papers. §\ref{sec:architecture} describes the
side-channel architecture and the three primitive tactics.
§\ref{sec:implementation} walks through the implementation of the §5
annotation and §6 augmentation+emission. §\ref{sec:evaluation} reports the
empirical comparison against Cyclist and presents worked examples.
§\ref{sec:limitations} identifies the boundary cases the current system does
not handle---most prominently merge-sort, whose recursion is not
structural---and outlines two extensions currently under separate
development: well-founded emission via Lee 2009-style measure
synthesis~\cite{lee2009ranking}, and Löb-induction emission via a
step-indexed embedding~\cite{atkey2013productive}.
